\section{同値双加群の障害消去とétale局所実現}\label{er:sec:etale}
本節では、可換模型の存在に使う局所実現原理を証明する。
入力は実在するsmooth properな代数的dg族であり、proper性だけの族ではない。
圏の全変形をHochschild複体から実現するという主張は用いず、
実在するperfect双加群のmoduliと、その完全な障害理論を用いる。
ここでのpairは自己同値一つを指定したもので、$\rho^3\simeq[2]$という
選択された同型やそのcoherenceを追加データとして含めない。

\subsection{perfect加群の小拡大障害}
小拡大は剰余体$k=\CC$を持つ局所Artin代数の全射
$A'\twoheadrightarrow A$で、核$I$が$\mathfrak m_{A'}I=0$を満たすものとする。
$D'$は$A'$上smooth properなdg代数とし、$D=D'\otL_{A'}A$とする。

\begin{lemma}[加群の完全な障害理論と特徴作用]\label{er:lem:obstruction}
$P\in\Perf(D)$の持上げには自然な障害
\[
 o(P;D')\in\operatorname{Ext}^2_{D_k}(P_k,P_k)\otimes_k I
 \tag{O}\label{er:eq:obstruction}
\]
があり、零であることと$\Perf(D')$への持上げの存在は同値である。
障害はMorita関手と対象の同型に関して自然である。
$D'=C'^{\rm op}\otL_{A'}B'$で$P$が$C$から$B$への同値を表すとき、
$B'$を固定し、$C'$のliftをHochschild類$\xi\in\mathrm{HH}^2(C_k)\otimes I$だけ変えると、
障害の変化は、符号を一貫して選べば
\[
 o_{\rm new}-o_{\rm old}=-\chi_P(\xi),\qquad
 \chi_P:\mathrm{HH}^2(C_k)\xrightarrow{\sim}
       \operatorname{Ext}^2_{C_k^{\rm op}\otimes B_k}(P_k,P_k).
 \tag{V}\label{er:eq:variation}
\]
\end{lemma}
\begin{proof}
Toën--Vaquié\cite[Theorem 3.6, Corollaries 3.15, 3.17]{TV}を
可換底$A'$と$D'$に適用する。smoothnessによりpseudo-perfect加群と
perfect加群は一致し、その実際の対象moduliの相対接複体は$\RHom_D(P,P)[1]$である。
導来対象moduliの平方零拡大に関する障害理論は、その接複体の次数一に完全な障害を持つ。
従って障害は$\operatorname{Ext}^2_D(P,P\otL_A I)$に入り、
導来随伴と$\mathfrak m_A I=0$によって\eqref{er:eq:obstruction}に一致する。
これは単なる接空間の計算を代数化の代用にしているのではない。
実在する対象の導来moduliが無限小凝集的であることを使っており、
障害のnull-homotopyはそのmoduliの$A'$-点、すなわち実際のperfect加群を与える。

変化式も説明する。cofibrantな代数模型と加群のbar模型を選び、
構造を次数一のbar余微分$b$で表す。任意のlift$\widetilde b$の曲率
$\widetilde b^{\,2}$は$I$に値を取り、$[b,\widetilde b^{\,2}]=0$である。
持上げを$I$値の$h$だけ変えると、$I^2=0$により曲率は$[b,h]$だけ変わる。
この曲率類が上のmoduliの障害である。
代数側を$\xi$だけ変えた際の曲率の変化は、$\xi$を加群の構造写像に挿入したcochainであり、
それが特徴作用$\chi_P(\xi)$である。符号はsourceの反対代数とtorsorの規約に合わせて
\eqref{er:eq:variation}と取る。必要なのは同じ規約を両式で使うことである。
この記述はLowen\cite[Theorem 4.8]{ERLowen}の特徴作用による障害の記述を、
bar模型でdg加群と任意の小拡大に書いたものである。
$P_k$の逆Morita双加群をテンソルすれば、その自己Homと対角の自己Homが同定されるので、
$\chi_P$は同型である。曲率の構成、またはmoduliの障害の自然性から、最後の自然性も従う。
\end{proof}

\begin{lemma}[還元後の同値とintertwiner]\label{er:lem:lift}
$C',B'$は$A'$上smooth properとする。
$P'\in\Perf(C'^{\rm op}\otL B')$の還元がMorita同値なら$P'$自身も同値である。
さらに$Q'$もperfectで、$\operatorname{Ext}^1(P_k,Q_k)=0$なら、
任意の射$P'_A\to Q'_A$は$A'$上へ持ち上がる。元の射が同型ならliftも同型となる。
\end{lemma}
\begin{proof}
perfect双加群は両側に随伴を持ち、smooth proper性により随伴双加群もperfectである。
unitとcounitのconeは底上perfectで、還元後に零であるから導来Nakayamaで零となる。
射の持上げには三角$Q'\otL I\to Q'\to Q'_A$を$\RHom(P',-)$で評価する。
射の障害は$\operatorname{Ext}^1(P_k,Q_k)\otimes I$に入り、仮定で零である。
持ち上げた射のconeに同じNakayamaを適用すれば同型性も従う。
\end{proof}

\subsection{rotationを保つ障害消去}
$H$は複素数体上smoothな有限型スキーム、$h\in H(\CC)$とする。
$\mathcal A_H$はsmooth proper dg族、$r_H$は可逆perfect双加群で表された自己同値とする。
$\mathcal A=\mathcal A_h$、$r=r_h$と置く。局所的なdg代数模型を取って議論する。

\begin{theorem}[自己同値付き局所実現の判定]\label{er:thm:criterion}
次を仮定する。
\[
 \mathrm{HH}^{<0}(\mathcal A)=0,\qquad
 \mathrm{HH}^{1}(\mathcal A)=0,
 \qquad
 T_hH\longrightarrow\mathrm{HH}^2(\mathcal A)^r
 \text{ が全射}.
 \tag{KS}\label{er:eq:KS}
\]
最後の写像は$\mathcal A_H$のKodaira--Spencer写像である。
$S$は有限型複素スキーム、$(\mathcal B_S,s_S)$はsmooth proper dg族と
可逆perfect双加群のpairとする。$s\in S(\CC)$に指定した同値
$(\mathcal A,r)\simeq(\mathcal B_s,s_s)$があるなら、
pointed étale近傍$V\to(S,s)$、$f:V\to H$、$f(v)=h$と
指定同値を延長するpairの同値
\[
 (f^*\mathcal A_H,f^*r_H)\simeq(\mathcal B_V,s_V)
\]
が存在する。
\end{theorem}
\begin{proof}
まず小拡大$A'\twoheadrightarrow A$上の持上げを証明する。
$\operatorname{Spec}A'\to S$を固定し、$A$上の点$h_A$、同値双加群$P_A$、
intertwiner $\alpha_A:s_AP_A\xrightarrow{\sim}P_Ar_A$が与えられているとする。
$H$はsmoothなので$h_A$の任意のlift$h'_0$を取れる。
両側のliftされた自己同値から、双加群圏の自己同値
\[
 \mathcal F'(P)=s'\circ P\circ(r'_0)^{-1}
\]
を得る。$\alpha_A$は$\mathcal F_A(P_A)\simeq P_A$を与える。
補題\ref{er:lem:obstruction}の自然性より、$P_A$の持上げ障害$o$は
$\mathcal F_k$の作用で不変である。
Morita同定$\operatorname{Ext}^2(P_k,P_k)=\mathrm{HH}^2(\mathcal A)$の下で、
この作用は$r$による共役作用（規約によりその逆）である。
従って
\[
 o\in\mathrm{HH}^2(\mathcal A)^r\otimes I.
 \tag{I}\label{er:eq:invariant}
\]
この段階では$\alpha_A$を$A'$上へ持ち上げたとは仮定していない。

\eqref{er:eq:KS}と\eqref{er:eq:variation}により、$h'_0$を
$T_hH\otimes I$の元だけ修正して$o$を消せる。
修正後のsourceと固定したtargetの間にperfect双加群$P'$が存在し、
補題\ref{er:lem:lift}によって同値となる。
$\alpha_A$を持ち上げる障害は
\[
 \operatorname{Ext}^1(s_kP_k,P_kr_k)\otimes I
       \simeq \mathrm{HH}^1(\mathcal A)\otimes I=0
\]
なので、intertwinerも同型として持ち上がる。
自己同値一つのintertwiner以外に群関係は課していないので、ここで必要なデータは全て揃う。
任意のnilpotent拡大は小拡大に分解できる。

次に代数化する。$H\times S$上で、perfectな
$\mathcal A_H^{\rm op}\otimes\mathcal B_S$-加群$P$と
同型$s_SP\simeq Pr_H$のmoduliを考え、$P$が同値である部分を$\mathcal E$とする。
Toën--Vaquié\cite[Theorem 3.6, Corollaries 3.15, 3.19]{TV}により、
perfect対象とその射の導来moduliは局所幾何的かつ局所有限表示である。
unit、counitおよびintertwinerのconeは底上perfectだから、その消滅は開条件である。
従って$\mathcal E$も局所幾何的かつ局所有限表示である。
同値点での自己Extの負次数は$\mathrm{HH}^{<0}(\mathcal A)=0$であり、
その近傍でこの消滅を保てる。従って古典的切断はその近傍でArtin 1-stackとなる。
ここでは局所dg代数模型を用いてTVの定理を$H\times S$上の可換係数環に適用している。

上の小拡大持上げは、指定点$e=(h,s,P_k,\alpha_k)$で
$\mathcal E\to S$が形式的にsmoothであることを示す。
局所有限表示のArtin stackの無限小判定により、この射は$e$でsmoothである。
$e$を覆うsmooth atlas $U\to\mathcal E$を取り、点の近傍へ縮めると$U\to S$もsmoothである。
滑らかな射のpointed étale局所section\cite[Tags 02HX, 055U]{StacksER}を使い、
$V\to S$と$V\to U\to\mathcal E$を得る。
この普遍対象が所要の族、双加群およびintertwinerを与える。
\end{proof}

\begin{remark}[使わない仮定と先行研究]\label{er:rem:etale-scope}
この証明は$\mathrm{HH}^3=0$、Hochschild cochainのformality、
抽象的なHochschild変形の全てがactual Morita変形になること、または$\rho^3=[2]$の指定を使わない。
Lin--Zhang\cite[Theorem 2.2, Proposition 2.4]{ERLinZhang}の同値核の障害消去と同じ原理を用いるが、
ここではtargetを非幾何的なsmooth proper dg族として扱い、
自己同値が障害を不変部分へ制限することと、intertwinerの持上げを別々に証明した。
同論文の幾何的projectorの仮定を満たさないtargetへその定理を直接適用していない。
この一般定式化の新規性は未確認である。
\end{remark}

\subsection{滑らかなcubicへの適用}
\begin{lemma}[低次Hochschild群と幾何的接写像]\label{er:lem:cubic}
$X=V(F)\subset\mathbb P^{N-1}$を滑らかなcubic、$N=n+2\ge5$とし、
$J(F)=\CC[z_1,\ldots,z_N]/(\partial_1F,\ldots,\partial_NF)$とする。
\[
 \mathrm{HH}^{<0}(\Ku(X))=0,\quad
 \mathrm{HH}^0(\Ku(X))=\CC,\quad \mathrm{HH}^1(\Ku(X))=0,
\]
\[
 \mathrm{HH}^2(\Ku(X))=
 \begin{cases}J(F)_3&n\ne4,\\J(F)_3\oplus\CC^2&n=4,\end{cases}
 \qquad
 H^1(X,T_X)\xrightarrow{\sim}\mathrm{HH}^2(\Ku(X))^{\rho_X}=J(F)_3.
 \tag{H2}\label{er:eq:HH}
\]
$\dim J(F)_3=\binom N3$である。
\end{lemma}
\begin{proof}
BFK\cite[Theorem 5.39]{BFKER}のequivariant matrix factorizationの公式を用いる。
$G=\mathbb G_m$、各座標の重みは1、potentialのcharacterは3である。
恒等sectorのKoszul cohomologyは、偏微分が正則列であるためJacobian商だけで、
cohomological degree $2r$に$J(F)_{3r}$を与える。
非恒等$g\in\mu_3$の固定空間は原点、conormalの次元は$N$、determinantの重みは$-N$である。
次数$q$の重み零条件は$(3q-N)/2=0$であり、従って$3\mid N$かつ$q=N/3$のときだけ
各sectorから一次元が現れる。これで負次数、次数0、1、2の群が得られる。

幾何的Kodaira--Spencer写像の同定は、単なる次元一致ではない。
$n\ge4$では$3\le(n+2)/2$なので、Huybrechts--Rennemo
\cite[Lemma 2.19, \S2.5, \S3.3]{HRER}のAtiyah--Kodaira--Spencerと
三次周期性の比較を適用でき、$F+\epsilon G$の像は恒等sectorの$[G]\in J(F)_3$である。
$n=3$ではJacovskis--Lin--Liu--Zhang\cite[Theorem 4.4(2)]{JLLZER}の
幾何的写像の同型を用いる。normal sequenceとEuler sequenceから$H^1(T_X)=J(F)_3$である。
幾何的族にはrotationがあるのでこの像は不変である。
$n=4$で余分な二方向が不変でないことは、Huybrechts\cite[Proposition 3.12]{HuyER}による。
rotationはsymplecticで、Mukai格子の$A_2$上で三つの根を巡回する。
不変なCalabi--Yau類による$\mathrm{HH}^2\simeq\mathrm{HH}_0$の同定の下で、
その二方向の固有値は$\zeta,\zeta^2$である。よって不変部分は正確に$J(F)_3$となる。
最後に偏微分が二次正則列であることから$H_{J(F)}(t)=(1+t)^N$を得る。
\end{proof}

\begin{theorem}[cubicのrotation付きétale局所実現]\label{er:thm:ER}
$S$は有限型複素スキーム、$\mathcal K_S$は実在するsmooth properな代数的dg族で、
$\rho_S$は可逆perfect双加群により表される自己同値とする。
$s\in S(\CC)$で、$n\ge3$の滑らかなcubic $X$に対する指定同値
$(\mathcal K_s,\rho_s)\simeq(\Ku(X),\rho_X)$があるなら、
pointed étale近傍$V\to(S,s)$、滑らかなcubic族$Y/V$、および指定同値を延長する同値
\[
 (\mathcal K_V,\rho_V)\simeq(\Ku(Y/V),\rho_{Y/V})
 \tag{ER}\label{er:eq:ER}
\]
が存在する。
\end{theorem}
\begin{proof}
滑らかな三次形式の射影パラメータ空間の$[F]$を含むアフィン開を$H$とする。
普遍cubic $\mathcal X/H$のrelative perfect圏はsmooth properで、
$\OO,\ldots,\OO(n-2)$のrelative exceptional列の射影により
smooth properな相対Kuznetsov成分を得る。
射影とrotationは有限個の評価三角で構成され、底変更に可換である。
普遍cubic上の有限個のtwistから生成子を選べるので、局所的なdg代数模型も存在する。
$T_{[F]}H\to H^1(T_X)=J(F)_3$は全射である。
補題\ref{er:lem:cubic}と定理\ref{er:thm:criterion}を適用する。
\end{proof}

\begin{corollary}[rotationを最初に指定しない存在]\label{er:cor:unpolarized}
$n\ne4$なら、前定理で$\rho_S$を入力として指定せずとも、
$\mathcal K_s\simeq\Ku(X)$からétale局所的な可換cubic模型が得られる。
これはrotationなしの大域的な一意性の主張ではない。
\end{corollary}
\begin{proof}
補題\ref{er:lem:cubic}では幾何的接写像が$\mathrm{HH}^2$全体に全射である。
定理\ref{er:thm:criterion}の証明から不変性とintertwinerの段階を省けばよい。
\end{proof}
